\documentclass[lettersize,journal]{IEEEtran}

\usepackage{amsmath,amssymb,amsfonts,bm}
\usepackage{algorithm}
\usepackage{algorithmic}
\usepackage{array}
\usepackage{booktabs}
\usepackage{cite}
\usepackage{graphicx}
\usepackage{makecell}
\usepackage{multirow}
\usepackage{url}
\usepackage{kotex}

\hyphenation{op-tical net-works semi-conduc-tor IEEE-Xplore}

\newcommand{\E}{\mathbb{E}}
\newcommand{\Var}{\operatorname{Var}}
\newcommand{\Cov}{\operatorname{Cov}}
\newcommand{\Bern}{\operatorname{Bernoulli}}
\newcommand{\N}{\mathcal{N}}
\newcommand{\PG}{\operatorname{PG}}
\newcommand{\IW}{\operatorname{IW}}
\newcommand{\R}{\mathbb{R}}
\newcommand{\D}{\mathcal{D}}
\newcommand{\diag}{\operatorname{diag}}
\newcommand{\sigmoid}{\sigma}
\newcommand{\given}{\,\middle|\,}

\begin{document}

\title{Hierarchical Fully Bayesian Approach
}
% Safe-Aware Personalized Adaptive Cruise Control by HOCBF-MPC?


\author{
    Jaewoong Han$^{1}$, Dongjae Kim$^{1}$, Taehyung Kim$^{2}$, Sheengil Kang$^{3}$, and Jongeun Choi$^{1}$
    \thanks{This work was supported by the Hyundai Motor Company, and also supported in part by the Institute of Information \& Communications Technology Planning \& Evaluation (IITP) grant funded by the Korea government (MSIT) (No. RS-2026-25524880)
    }% <-this % stops a space
    \thanks{$^{1}$Jaewoong Han, Dongjae Kim, and Jongeun Choi are with the School of Mechanical Engineering, Yonsei University, Seoul, South Korea.
    {\tt\small \{hanjw2496, kimdj98, jongeunchoi\}@yonsei.ac.kr}}%
    \thanks{$^{2}$Taehyung Kim is with the Department of Mobility Systems Engineering, Yonsei University, Seoul, South Korea.
    {\tt\small kth8606@yonsei.ac.kr}}
    \thanks{$^{3}$Sheengil Kang is with SDV Division, Hyundai Motors, Hwaseong-si, South Korea.
    {\tt\small sheengil@hyundai.com}}
}
\maketitle

\begin{abstract}
\end{abstract}

\begin{IEEEkeywords}
\end{IEEEkeywords}


\section{Introduction}

\section{Related Work}

\section{Preliminaries}

\section{Methodology}

Each item corresponds to a driving segment or trajectory window $\tau_{ui}$.
We first map $\tau_{ui}$ to a feature vector through a deterministic feature extractor:
\begin{equation*}
    \bm z_{ui}=\phi(\tau_{ui})\in\mathbb{R}^{d}.
\end{equation*}
For user $u$, we define the latent personalized reward function as
\begin{equation*}
    r_{ui} = r_u(\tau_{ui})
    = \bm\theta_u^\top \bm z_{ui}
    = \bm\theta_u^\top \phi(\tau_{ui}).
\end{equation*}

The binary multi-user feedback $y_{ui} \in \{0, 1\}$ is labeled based on this personal reward. To express this observation process using an explicit equality form while maintaining the logistic choice structure, we introduce a latent variable formulation:
\begin{equation*}
y_{ui} = \mathbb{I}(r_{ui} + \eta_{ui} > 0), \qquad
\eta_{ui}\sim \mathrm{Logistic}(0,1).
\end{equation*}
where $\mathbb{I}(\cdot)$ is the indicator function, and $\eta_{ui}$ represents the random observation noise following a standard logistic distribution. By integrating out the latent logistic noise $\eta_{ui}$, this formulation can be equivalently expressed as an explicit choice probability form:
\begin{equation*}
P(y_{ui} = 1 \mid \bm \theta _u, \bm z_{ui}) = \sigmoid \! \left( \bm \theta _u^{\top } \bm z_{ui } \right) .
\end{equation*}

Consequently, the generative process for the binary feedback is modeled as a Bernoulli trial:
\begin{equation*}
    y_{ui}\mid\bm\theta_u,\bm z_{ui}
    \sim\Bern\!\left(
        \sigmoid(\bm\theta_u^{\top}\bm z_{ui})
    \right).
\end{equation*}

Users are heterogeneous but not unrelated. We model user-specific weights as draws from a shared population distribution,
\begin{equation*}
    \bm\theta_u\mid\bm\mu,\Sigma
    \sim\N(\bm\mu,\Sigma),
    \qquad u=1,\ldots,U.
\end{equation*}
Here, $\bm\mu$ represents the population-average preference vector, while $\Sigma$ represents variation and correlation in preference dimensions across users.

We place a Normal--Inverse-Wishart prior on the population parameters:
\begin{align}
    \Sigma &\sim \IW(\nu_0,\Lambda_0), \\
    \bm\mu\mid\Sigma
    &\sim
    \N\!\left(\bm m_0,\frac{1}{\kappa_0}\Sigma\right).
    \label{eq:niw_prior}
\end{align}

The target posterior is
\begin{equation*}
    p(\bm\mu,\Sigma,\bm\theta_{1:U}
    \mid \D_{1:U}),
\end{equation*}
where $\bm\theta_{1:U}=\{\bm\theta_1,\ldots,\bm\theta_U\}$.


\subsection{User-Specific Reward Weight}
For each user $u$, let
\begin{equation*}
Z_u = \begin{bmatrix} 1\\ \bm z_{u1}^{\top}\\ \vdots\\ \bm z_{uN_u}^{\top} \end{bmatrix}
\in\mathbb{R}^{N_u\times d}, \qquad \bm y_u=(y_{u1},\ldots,y_{uN_u})^\top .
\end{equation*}

The Bernoulli-logistic likelihood is not conjugate to the Gaussian prior. 
To obtain conditionally Gaussian updates for $\bm\theta_u$, we use the Pólya--Gamma augmentation. 
Let $\psi_{ui}=\bm\theta_u^\top\bm z_{ui}$ and $c_{ui}=y_{ui}-1/2$. 
Then the likelihood can be written as
\begin{align}
p(y_{ui}\mid \psi_{ui})
&\propto \exp(c_{ui}\psi_{ui}) \nonumber \\
&\quad \int_0^\infty \exp\!\left(-\frac{\omega_{ui}\psi_{ui}^2}{2}\right) p_{\PG}(\omega_{ui}\mid 1,0) \,d\omega_{ui}.
\end{align}

Conditioned on the current user weights,
\begin{equation*}
    \omega_{ui}\mid\bm\theta_u,\D_u
    \sim\PG\!\left(1,\bm\theta_u^\top \bm z_{ui}\right).
\end{equation*}

For user $u$, define
\begin{equation*}
    \Omega_u=\diag(\omega_{u1},\ldots,\omega_{uN_u}),
    \qquad
    \bm c_u=\bm y_u-\frac{1}{2}\bm 1_{N_u}.
\end{equation*}
Then
\begin{equation*}
    \bm\theta_u\mid
    \D_u,\bm\omega_u,\bm\mu,\Sigma
    \sim\N(\bm m_u,V_u),
\end{equation*}
with
\begin{align}
    V_u
    &=\left(Z_u^{\top}\Omega_u Z_u+\Sigma^{-1}\right)^{-1},
    \label{eq:user_cov}\\
    \bm m_u
    &=V_u\left(Z_u^{\top}\bm c_u+\Sigma^{-1}\bm\mu\right).
    \label{eq:user_mean}
\end{align}


\subsection{Population Parameters}
Conditioned on the current draws $\bm\theta_1,\ldots,\bm\theta_U$, the Normal--Inverse-Wishart posterior is available in closed form \cite{murphy2007}. Let
\begin{equation*}
    \bar{\bm\theta}
    =\frac{1}{U}\sum_{u=1}^{U}\bm\theta_u
\end{equation*}
and define
\begin{align}
    \kappa_U&=\kappa_0+U,
    \label{eq:kappa_update}\\
    \nu_U&=\nu_0+U,
    \notag\\
    \bm m_U
    &=\frac{\kappa_0\bm m_0+U\bar{\bm\theta}}
    {\kappa_0+U},
    \notag
\end{align}
and
\begin{equation}
\begin{split}
    \Lambda_U
    ={}&\Lambda_0
    +\sum_{u=1}^{U}
    (\bm\theta_u-\bar{\bm\theta})
    (\bm\theta_u-\bar{\bm\theta})^{\top}\\
    &+\frac{\kappa_0U}{\kappa_0+U}
    (\bar{\bm\theta}-\bm m_0)
    (\bar{\bm\theta}-\bm m_0)^{\top}.
    \label{eq:lambda_update}
\end{split}
\end{equation}
The conditional distributions are
\begin{align}
    \Sigma\mid\bm\theta_{1:U}
    &\sim\IW(\nu_U,\Lambda_U),
    \label{eq:sigma_conditional}\\
    \bm\mu\mid\Sigma,\bm\theta_{1:U}
    &\sim\N\!\left(
        \bm m_U,\frac{1}{\kappa_U}\Sigma
    \right).
    \label{eq:mu_conditional}
\end{align}

\subsection{New-User Personalization}
For a new user $*$, the learned population posterior provides a personalized prior. 
For each posterior sample $(\bm\mu^{(m)},\Sigma^{(m)})$, we initialize
\begin{equation*}
    \bm\theta_*^{(m)}
    \sim
    \N(\bm\mu^{(m)},\Sigma^{(m)}).
\end{equation*}
Given a small context set
\begin{equation*}
    \D_*^{\mathrm{ctx}}
    =
    \{(\bm z_{*i},y_{*i})\}_{i=1}^{N_*},
\end{equation*}
we update $\bm\theta_*^{(m)}$ using the same Pólya--Gamma conditional Gaussian update:
\begin{align}
    V_*^{(m)}
    &=
    \left(
    Z_*^\top \Omega_*^{(m)} Z_*
    +
    (\Sigma^{(m)})^{-1}
    \right)^{-1},\\
    \bm m_*^{(m)}
    &=
    V_*^{(m)}
    \left(
    Z_*^\top \bm c_*
    +
    (\Sigma^{(m)})^{-1}\bm\mu^{(m)}
    \right),
\end{align}
\begin{equation*}
    \bm\theta_*^{(m)}
    \sim
    \N(\bm m_*^{(m)},V_*^{(m)}).
\end{equation*}

For a query item with feature vector $\bm z$, the posterior predictive probability is approximated by Monte Carlo integration:
\begin{equation}
    P(y=1\mid \bm z,\D_*^{\mathrm{ctx}},\D_{1:U})
    \approx
    \frac{1}{M}
    \sum_{m=1}^{M}
    \sigma\!\left((\bm\theta_*^{(m)})^\top\bm z\right).
    \label{eq:posterior_predictive}
\end{equation}

The epistemic and aleatoric uncertainty of the prediction can be estimated by the posterior variance
\begin{equation}
\underbrace{\operatorname{Var}(Y\mid z,\mathcal D)}_{\text{total predictive uncertainty}}
= 
\underbrace{\mathbb E_m[p^{(m)}(1-p^{(m)})]}_{\text{aleatoric}}
+
\underbrace{\operatorname{Var}_m[p^{(m)}]}_{\text{epistemic}}
\end{equation}
where $p^{(m)}=\sigma((\theta_*^{(m)})^\top z)$

\section{Next Progress}
\subsection{Feature Selection}
\label{sec:feature_selection}

The current framework assumes that the feature extractor
$\phi(\tau_{ui})$ is predefined. However, including all available
sensor-derived features may introduce redundant information, increase
computational cost, and reduce interpretability. We therefore consider
two complementary feature-selection criteria: (i) projection-predictive
selection based on the Kullback--Leibler (KL) divergence from the full
Bayesian model, and (ii) conditional predictive information gain based
on the observed binary feedback.

The two criteria address different questions. The projection-predictive
criterion identifies a compact subset that preserves the predictions of
the full hierarchical model, whereas the conditional information-gain
criterion identifies features that provide additional information about
the observed feedback after conditioning on the features already
selected. Neither criterion performs active query selection or attempts
to reduce new-user parameter uncertainty.

\subsubsection{Selection Units}

Let $\mathcal{J}=\{1,\ldots,d\}$ denote the set of individual feature
indices and let $\mathcal{G}=\{1,\ldots,G\}$ denote the set of sensors.
For sensor $g$, define
\begin{equation}
    \mathcal{J}_g
    =
    \left\{
        j\in\mathcal{J}:
        \text{feature }j\text{ is derived from sensor }g
    \right\}.
\end{equation}
Thus, $\mathcal{J}_g$ may contain several trajectory-level statistics,
such as the mean, variance, maximum, or temporal variation of the same
sensor signal. The complete feature vector can be written as
\begin{equation}
    \bm z_{ui}
    =
    \left[
        \bm z_{ui,\mathcal{J}_1}^{\top},
        \ldots,
        \bm z_{ui,\mathcal{J}_G}^{\top}
    \right]^{\top}.
\end{equation}

We consider the following two sets of selectable units:
\begin{align}
    \mathcal{A}_{\mathrm{sensor}}
    &=
    \left\{
        \mathcal{J}_1,\ldots,\mathcal{J}_G
    \right\},\\
    \mathcal{A}_{\mathrm{feature}}
    &=
    \left\{
        \{1\},\ldots,\{d\}
    \right\}.
\end{align}
Selecting $A=\mathcal{J}_g$ therefore selects all statistics extracted
from sensor $g$, whereas selecting $A=\{j\}$ selects one individual
sensor--statistic feature. The two selection procedures described below
are applied separately at the sensor and individual-feature levels.

Let $S\subseteq\mathcal{J}$ denote the set of currently selected feature
indices, and let $\bm z_{ui,S}$ denote the corresponding feature vector.
The bias column, when included, is always retained and is not treated as
a selectable unit. Its index set is denoted by $S_0$.

\subsubsection{Projection-Predictive Selection}
\label{sec:projection_selection}

The full hierarchical Bayesian model is first fitted using all
available features. For posterior draw $m$, the reference predictive
probability for observation $(u,i)$ is
\begin{equation}
    p_{ui}^{(m)}
    =
    \sigmoid\left(
        (\bm\theta_u^{(m)})^{\top}\bm z_{ui}
    \right).
    \label{eq:reference_probability}
\end{equation}
These probabilities contain the predictive behavior learned by the full
model, including both the population-level structure and
user-specific preferences.

For a candidate subset $S$, we project each reference posterior draw
onto a restricted logistic model using only $\bm z_{ui,S}$. The
projected coefficient is defined as
\begin{equation}
\begin{split}
\widetilde{\bm\theta}_{u,S}^{(m)} &= \arg\min_{\bm\beta\in\mathbb{R}^{|S|}} \Bigg\{ \frac{1}{N_u} \sum_{i=1}^{N_u} D_{\mathrm{KL}} \Big[ \Bern \big(p_{ui}^{(m)} \big) \\
&\quad \Vert \Bern \big(\sigmoid (\bm\beta^{\top} \bm z_{ui,S}) \big) \Big] + \frac{\lambda_{\mathrm{proj}}}{2} \left\| \bm\beta_{\mathrm{var}} \right\|_2^2 \Bigg\},
\label{eq:projection_parameter}
\end{split}
\end{equation}

where $\bm\beta_{\mathrm{var}}$ contains the coefficients associated
with nonconstant feature columns. Constant columns, including the bias,
are not penalized.

The projected predictive probability is then
\begin{equation}
    q_{ui,S}^{(m)}
    =
    \sigmoid\left(
        (\widetilde{\bm\theta}_{u,S}^{(m)})^{\top}
        \bm z_{ui,S}
    \right).
\end{equation}
For Bernoulli distributions, the pointwise KL divergence is
\begin{equation}
\begin{split}
    &D_{\mathrm{KL}} \left[ \Bern(p)\Vert\Bern(q) \right] \\
    &\quad = p\log\frac{p}{q} + (1-p)\log\frac{1-p}{1-q}.
\label{eq:bernoulli_kl}
\end{split}
\end{equation}

Averaging over users, observations, and a set
$\mathcal{M}_{\mathrm{proj}}$ of retained posterior draws gives the
projection loss
\begin{equation}
\begin{split}
\mathcal{K}(S) &= \frac{1}{|\mathcal{M}_{\mathrm{proj}}| \sum_{u=1}^{U} N_u} \sum_{m \in \mathcal{M}_{\mathrm{proj}}} \\
&\quad \sum_{u=1}^{U} \sum_{i=1}^{N_u} D_{\mathrm{KL}} \left [\Bern\big(p_{ui}^{(m)}\big) \Vert \Bern\big(q_{ui,S}^{(m)}\big) \right ] .
\label{eq:projection_loss}
\end{split}
\end{equation}

A smaller value of $\mathcal{K}(S)$ indicates that the restricted
feature set more faithfully reproduces the predictions of the full
Bayesian model.

Starting from the bias-only model, the next selectable unit is chosen by
greedy forward selection:
\begin{equation}
    A_t^{*}
    =
    \arg\min_{
        A\in\mathcal{A},
        \,A\cap S_t=\emptyset
    }
    \mathcal{K}(S_t\cup A),
    \label{eq:projection_forward}
\end{equation}
followed by
\begin{equation}
    S_{t+1}=S_t\cup A_t^{*},
    \qquad S_0=\{\text{bias}\}.
\end{equation}

To define a stopping criterion, we measure the fraction of the
bias-model projection loss removed by subset $S$:
\begin{equation}
    R_{\mathrm{KL}}(S)
    =
    1-
    \frac{\mathcal{K}(S)}
         {\mathcal{K}(S_0)}.
    \label{eq:kl_captured}
\end{equation}
The smallest subset satisfying
\begin{equation}
    R_{\mathrm{KL}}(S)\geq \rho
\end{equation}
is selected, where $\rho$ denotes the desired predictive-information
retention level. In the current implementation, $\rho=0.95$ is used.

The projection in \eqref{eq:projection_parameter} is a deterministic
soft-label logistic optimization. Consequently, Gibbs sampling is not
repeated for every candidate subset. The observed feedback does not
appear directly in the projection objective, but it affects the
selection indirectly through the full-model posterior samples used in
\eqref{eq:projection_forward}.

\subsubsection{Conditional Predictive Information Gain}
\label{sec:conditional_information_selection}

Projection-predictive selection measures agreement with the full model,
but it does not directly quantify how much additional information a
feature provides about the observed feedback. We therefore introduce a
second criterion based on conditional mutual information.

Let $\mathsf{U}$ denote user identity. For a candidate unit $A$ and a
currently selected feature set $S$, we define
\begin{equation}
\begin{split}
    \mathcal{I}(A\mid S)
    &=
    I\left(
        Y;\bm Z_A
        \mid
        \bm Z_S,\mathsf{U}
    \right)\\
    &=
    H\left(
        Y\mid\bm Z_S,\mathsf{U}
    \right)
    -
    H\left(
        Y\mid\bm Z_{S\cup A},\mathsf{U}
    \right).
    \label{eq:conditional_mi}
\end{split}
\end{equation}
Equivalently,
\begin{equation}
\begin{split}
    \mathcal{I}(A\mid S)
    =
    \E
    \left[
        \log
        \frac{
            p(Y\mid\bm Z_{S\cup A},\mathsf{U})
        }{
            p(Y\mid\bm Z_S,\mathsf{U})
        }
    \right].
    \label{eq:conditional_mi_logratio}
\end{split}
\end{equation}
Thus, $\mathcal{I}(A\mid S)$ measures the additional predictive
information provided by $A$ after accounting for the selected features
and user identity. Conditioning on $\mathsf{U}$ prevents between-user
differences from being incorrectly attributed to sensor features.

The conditional probabilities in
\eqref{eq:conditional_mi_logratio} are estimated using regularized
user-conditioned logistic models. For subset $S$, the auxiliary
conditional model has the form
\begin{equation}
    \widehat p_S
    \left(
        Y_{ui}=1\mid\bm z_{ui,S},u
    \right)
    =
    \sigmoid\left(
        \alpha_u+
        \bm\beta_{u,S}^{\top}\bm z_{ui,S}
    \right),
    \label{eq:user_conditioned_auxiliary}
\end{equation}
where each user has a separate intercept and separate feature
coefficients. The coefficients are estimated with an $\ell_2$
regularization term.

To reduce in-sample optimism, the information gain is evaluated using
within-user stratified cross-fitting. Let $k(ui)$ denote the validation
fold containing observation $(u,i)$, and let
$\widehat p_{S}^{(-k(ui))}$ denote a conditional model fitted without
that fold. Define the binary log score as
\begin{equation}
    \ell(y,p)
    =
    y\log p+(1-y)\log(1-p).
\end{equation}
The user-specific cross-fitted information gain is
\begin{equation}
\begin{split}
    \widehat{\mathcal{I}}_u(A\mid S)
    =
    \frac{1}{N_u}
    \sum_{i=1}^{N_u}
    \Big[
        &\ell\left(
            y_{ui},
            \widehat p_{S\cup A}^{(-k(ui))}
        \right)
        \nonumber\\
        &-
        \ell\left(
            y_{ui},
            \widehat p_S^{(-k(ui))}
        \right)
    \Big].
    \label{eq:user_cmi_estimate}
\end{split}
\end{equation}
We assign equal weight to each user and define the overall estimate as
\begin{equation}
    \widehat{\mathcal{I}}(A\mid S)
    =
    \frac{1}{U}
    \sum_{u=1}^{U}
    \widehat{\mathcal{I}}_u(A\mid S).
    \label{eq:cmi_user_average}
\end{equation}
Its user-level standard error is estimated by
\begin{equation}
\begin{split}
    \widehat{\operatorname{SE}}(A\mid S)
    =
    \left[
        \frac{1}{U(U-1)}
        \sum_{u=1}^{U}
        \left(
            \widehat{\mathcal{I}}_u(A\mid S)
            -
            \widehat{\mathcal{I}}(A\mid S)
        \right)^2
    \right]^{1/2}.
    \label{eq:cmi_user_se}
\end{split}
\end{equation}

Starting from the user-intercept-only model, the next selectable unit is
chosen as
\begin{equation}
    A_t^{*}
    =
    \arg\max_{
        A\in\mathcal{A},
        \,A\cap S_t=\emptyset
    }
    \widehat{\mathcal{I}}(A\mid S_t),
    \label{eq:cmi_forward_selection}
\end{equation}
and the selected set is updated according to
\begin{equation}
    S_{t+1}=S_t\cup A_t^{*}.
\end{equation}
The forward search terminates when the best remaining candidate has
nonpositive information gain,
\begin{equation}
    \max_A
    \widehat{\mathcal{I}}(A\mid S_t)
    \leq \epsilon_{\mathrm{IG}},
    \label{eq:cmi_stopping}
\end{equation}
where $\epsilon_{\mathrm{IG}}=0$ in the current implementation. A more
conservative indicator of feature importance is
\begin{equation}
    \widehat{\mathcal{I}}(A\mid S)
    -
    \widehat{\operatorname{SE}}(A\mid S)
    >
    \epsilon_{\mathrm{IG}}.
    \label{eq:cmi_importance}
\end{equation}

Unlike projection-predictive selection, the conditional information-gain
criterion directly uses the observed feedback $y_{ui}$ during both
conditional-model fitting and held-out log-score evaluation. It should
therefore be interpreted as a supervised, model-based estimate of
conditional predictive information. Equality with the exact
conditional mutual information in \eqref{eq:conditional_mi} requires
correctly specified conditional probability models.

\subsubsection{Hierarchical Spike-and-Slab Selection}
\label{sec:spike_slab_selection}

The preceding selection criteria are applied after fitting a reference
model or through an auxiliary predictive model. As an alternative, we
incorporate feature selection directly into the hierarchical Bayesian
model using binary inclusion variables.

Let $\mathcal{A}=\{A_1,\ldots,A_K\}$ denote the selectable units defined
previously. For individual-feature selection, $A_g$ contains one
feature index, whereas for sensor-level selection, $A_g$ contains all
statistics derived from sensor $g$. We associate each unit with a
global binary inclusion indicator
\begin{equation}
    \gamma_g\in\{0,1\},
    \qquad g=1,\ldots,K.
\end{equation}
The same $\gamma_g$ is shared across users, so the resulting selection
represents population-level relevance rather than user-specific
selection.

Let $\widetilde{\bm\gamma}\in\{0,1\}^{d}$ denote the feature-level mask
obtained by expanding the unit-level indicators:
\begin{equation}
    \widetilde{\gamma}_j
    =
    \gamma_g,
    \qquad j\in A_g.
\end{equation}
When a bias term is included, its indicator is fixed as
$\widetilde{\gamma}_{\mathrm{bias}}=1$. Define
\begin{equation}
    D_{\bm\gamma}
    =
    \diag(\widetilde{\bm\gamma}).
\end{equation}
The effective user coefficient is then
\begin{equation}
    \bm\beta_u
    =
    D_{\bm\gamma}\bm\theta_u,
    \label{eq:effective_spike_slab_coefficient}
\end{equation}
where $\bm\theta_u$ denotes the latent slab coefficient satisfying the
original hierarchical prior
\begin{equation}
    \bm\theta_u\mid\bm\mu,\Sigma
    \sim
    \N(\bm\mu,\Sigma).
\end{equation}
Consequently, the binary-feedback model becomes
\begin{equation}
    y_{ui}
    \mid
    \bm\theta_u,\bm\gamma,\bm z_{ui}
    \sim
    \Bern\left(
        \sigmoid\left(
            \bm z_{ui}^{\top}
            D_{\bm\gamma}\bm\theta_u
        \right)
    \right).
    \label{eq:spike_slab_likelihood}
\end{equation}
If $\gamma_g=0$, every effective coefficient associated with unit
$A_g$ is exactly zero for all users. If $\gamma_g=1$, the corresponding
coefficients follow the hierarchical Gaussian slab distribution.

We place a Beta--Bernoulli prior on the inclusion indicators:
\begin{align}
    \pi
    &\sim
    \operatorname{Beta}(a_{\pi},b_{\pi}),\\
    \gamma_g\mid\pi
    &\sim
    \Bern(\pi),
    \qquad g=1,\ldots,K.
    \label{eq:gamma_prior}
\end{align}
The parameter $\pi$ represents the population-level proportion of
included units. Its conditional posterior is
\begin{equation}
    \pi\mid\bm\gamma
    \sim
    \operatorname{Beta}\left(
        a_{\pi}+\sum_{g=1}^{K}\gamma_g,\,
        b_{\pi}+K-\sum_{g=1}^{K}\gamma_g
    \right).
    \label{eq:pi_conditional}
\end{equation}

The Pólya--Gamma augmentation permits direct Gibbs updates of each
inclusion indicator. Let
\begin{equation}
    c_{ui,g}
    =
    \bm z_{ui,A_g}^{\top}
    \bm\theta_{u,A_g}
\end{equation}
denote the contribution of candidate unit $A_g$, and define the linear
predictor excluding this unit as
\begin{equation}
    \eta_{ui,-g}
    =
    \bm z_{ui}^{\top}
    D_{\bm\gamma_{-g}}\bm\theta_u
    +
    \epsilon_{ui}.
\end{equation}
Here, $\epsilon_{ui}=0$ when the additional observation-noise term is
disabled. Since the augmented likelihood is proportional to
\begin{equation}
    \exp\left(
        c_{ui}\psi_{ui}
        -
        \frac{\omega_{ui}}{2}\psi_{ui}^{2}
    \right),
    \qquad
    c_{ui}=y_{ui}-\frac{1}{2},
\end{equation}
the conditional log-posterior odds of including unit $A_g$ are
\begin{equation}
\begin{split}
    \Delta_g
    ={}&
    \log\frac{\pi}{1-\pi}\\
    &+
    \sum_{u=1}^{U}
    \sum_{i=1}^{N_u}
    \left[
        c_{ui}c_{ui,g}
        -
        \frac{\omega_{ui}}{2}
        \left\{
            (\eta_{ui,-g}+c_{ui,g})^2
            -
            \eta_{ui,-g}^2
        \right\}
    \right].
    \label{eq:gamma_log_odds}
\end{split}
\end{equation}
Accordingly, the full conditional distribution is
\begin{equation}
    \gamma_g\mid\mathrm{rest}
    \sim
    \Bern\left(
        \sigmoid(\Delta_g)
    \right).
    \label{eq:gamma_conditional}
\end{equation}
The indicators are updated sequentially within each Gibbs iteration,
conditioning on the most recent values of the remaining indicators.

Given $\bm\gamma$, define the masked design matrix
\begin{equation}
    Z_{u,\bm\gamma}
    =
    Z_uD_{\bm\gamma}.
\end{equation}
The conditional posterior of the slab coefficient remains Gaussian:
\begin{equation}
    \bm\theta_u
    \mid
    \bm\omega_u,\bm\gamma,\bm\mu,\Sigma,\mathcal D_u
    \sim
    \N\left(
        \bm m_{u,\bm\gamma},
        V_{u,\bm\gamma}
    \right),
\end{equation}
where
\begin{align}
    V_{u,\bm\gamma}
    &=
    \left(
        Z_{u,\bm\gamma}^{\top}
        \Omega_u
        Z_{u,\bm\gamma}
        +
        \Sigma^{-1}
    \right)^{-1},\\
    \bm m_{u,\bm\gamma}
    &=
    V_{u,\bm\gamma}
    \left[
        Z_{u,\bm\gamma}^{\top}
        \left(
            \bm c_u-\Omega_u\bm\epsilon_u
        \right)
        +
        \Sigma^{-1}\bm\mu
    \right].
    \label{eq:masked_theta_conditional}
\end{align}
Thus, introducing $\bm\gamma$ preserves the conditionally conjugate
Gaussian update for each user. The Normal--Inverse-Wishart updates of
$\bm\mu$ and $\Sigma$ are unchanged and are applied to the latent slab
coefficients $\bm\theta_{1:U}$.

For retained posterior draw $m$, the effective coefficient used for
prediction is
\begin{equation}
    \bm\beta_u^{(m)}
    =
    D_{\bm\gamma^{(m)}}
    \bm\theta_u^{(m)}.
\end{equation}
For a new user, we first draw
\begin{equation}
    \bm\theta_*^{(m)}
    \sim
    \N\left(
        \bm\mu^{(m)},\Sigma^{(m)}
    \right)
\end{equation}
and use
\begin{equation}
    \bm\beta_*^{(m)}
    =
    D_{\bm\gamma^{(m)}}
    \bm\theta_*^{(m)}.
\end{equation}
During new-user personalization, the population-level indicator
$\bm\gamma^{(m)}$ remains fixed, while the corresponding slab
coefficient $\bm\theta_*^{(m)}$ is updated using the observed context.

The posterior inclusion probability of unit $A_g$ is estimated as
\begin{equation}
    \operatorname{PIP}_g
    =
    P(\gamma_g=1\mid\mathcal D)
    \approx
    \frac{1}{M}
    \sum_{m=1}^{M}
    \gamma_g^{(m)}.
    \label{eq:posterior_inclusion_probability}
\end{equation}
A unit with a high $\operatorname{PIP}_g$ is frequently included in
posterior models. The subset
\begin{equation}
    \widehat{\mathcal A}_{0.5}
    =
    \left\{
        A_g:
        \operatorname{PIP}_g\geq 0.5
    \right\}
\end{equation}
is reported as the median-probability model.

This method directly uses the observed feedback through the joint
posterior and differs from the projection-predictive and conditional
information-gain criteria. It modifies the hierarchical generative
model rather than applying a separate post-hoc ranking criterion.
However, correlated features may divide their posterior inclusion
probabilities, and the binary indicators can exhibit slow switching.
Therefore, posterior inclusion probabilities should be interpreted
together with the indicator traces and the number of posterior
transitions. Moreover, because $\bm\gamma$ is shared across users, the
current formulation may suppress a feature that is important only for
a small subset of users.



\subsubsection{Interpretation and Limitations}

The two criteria provide complementary interpretations. A low
$\mathcal{K}(S)$ indicates that subset $S$ preserves the predictive
behavior of the full hierarchical Bayesian model, while a large
$\widehat{\mathcal{I}}(A\mid S)$ indicates that candidate $A$ provides
additional information about the observed feedback beyond the features
already selected.

Both methods use greedy forward selection and therefore do not guarantee
a globally optimal subset. In addition, highly correlated features may
substitute for one another, so a feature assigned low incremental
importance is not necessarily uninformative in isolation. At the sensor
level, groups containing more derived statistics also have greater
representational capacity than smaller groups. Finally, the conditional
information-gain estimator evaluates within-user predictive relevance;
it does not evaluate generalization to an entirely unseen user.

Accordingly, the projection-predictive and conditional-information
procedures are reported separately as two feature-selection analyses.
The resulting subsets should be treated as candidate reduced
representations and subsequently verified by refitting and evaluating
the hierarchical model on held-out data.


\subsection{Add Bias Term}

\subsection{Scalable Bayesian Reward Learning}


\section{Experiments}

\section{Conclusion and Future Work}



\begin{algorithm}[t]
\caption{Population-Level Bayesian Reward Learning}
\label{alg:population_learning}
\begin{algorithmic}[1]
\REQUIRE Multi-user data $\D_{1:U}$, NIW hyperparameters, burn-in $B$, retained samples $M$
\ENSURE Posterior samples $\{ \bm\mu^{(m)},\Sigma^{(m)},\bm\theta_{1:U}^{(m)} \}_{m=1}^{M}$

\STATE Extract and standardize features $\bm z_{ui}=\phi(\tau_{ui})$ for all users.
\STATE Initialize $\bm\mu$, $\Sigma$, and $\{\bm\theta_u\}_{u=1}^{U}$.
\STATE Set $s\leftarrow 0$.

\FOR{$t=1$ to $B+M$}
    \FOR{$u=1$ to $U$}
        \STATE Sample $\omega_{ui}\sim\PG(1,\bm\theta_u^\top\bm z_{ui})$ for $i=1,\ldots,N_u$.
        \STATE Sample $\bm\theta_u\mid\D_u,\bm\omega_u,\bm\mu,\Sigma$ using \eqref{eq:user_cov}--\eqref{eq:user_mean}.
    \ENDFOR

    \STATE Compute NIW posterior parameters using \eqref{eq:kappa_update}--\eqref{eq:lambda_update}.
    \STATE Sample $\Sigma\mid\bm\theta_{1:U}$ using \eqref{eq:sigma_conditional}.
    \STATE Sample $\bm\mu\mid\Sigma,\bm\theta_{1:U}$ using \eqref{eq:mu_conditional}.

\ENDFOR

\RETURN $\{ \bm\mu^{(m)},\Sigma^{(m)},\bm\theta_{1:U}^{(m)} \}_{m=1}^{M}$
\end{algorithmic}
\end{algorithm}




\begin{algorithm}[t]
\caption{New-User Personalization and Prediction}
\label{alg:new_user_prediction}
\begin{algorithmic}[1]
\REQUIRE Population posterior samples $\{ \bm\mu^{(m)},\Sigma^{(m)} \}_{m=1}^{M}$, context data $\D_*^{\mathrm{ctx}}$, query features $\{\bm z_j\}_{j=1}^{N_q}$, adaptation sweeps $L$
\ENSURE Predictive probabilities $\{\hat p_j\}_{j=1}^{N_q}$

\FOR{$m=1$ to $M$}
    \STATE Initialize $\bm\theta_*^{(m)}\sim\N(\bm\mu^{(m)},\Sigma^{(m)})$.
\ENDFOR

\FOR{$\ell=1$ to $L$}
    \FOR{$m=1$ to $M$}
        \STATE Sample PG variables $\omega_{*i}^{(m)}\sim\PG(1,(\bm\theta_*^{(m)})^\top\bm z_{*i})$.
        \STATE Update $\bm\theta_*^{(m)}$ using the conditional Gaussian update in the new-user equations.
    \ENDFOR
\ENDFOR

\FOR{$j=1$ to $N_q$}
    \STATE Compute $\hat p_j$ using \eqref{eq:posterior_predictive}.
    \STATE Estimate epistemic uncertainty by the posterior variance of $\{\sigma((\bm\theta_*^{(m)})^\top\bm z_j)\}_{m=1}^{M}$.
\ENDFOR

\RETURN $\{\hat p_j\}_{j=1}^{N_q}$
\end{algorithmic}
\end{algorithm}








\bibliographystyle{IEEEtran}
\bibliography{refs}
\end{document}
